% Options for packages loaded elsewhere
\PassOptionsToPackage{unicode}{hyperref}
\PassOptionsToPackage{hyphens}{url}
%
\documentclass[
  11pt,
]{article}
\usepackage{amsmath,amssymb}
\usepackage{iftex}
\ifPDFTeX
  \usepackage[T1]{fontenc}
  \usepackage[utf8]{inputenc}
  \usepackage{textcomp} % provide euro and other symbols
\else % if luatex or xetex
  \usepackage{unicode-math} % this also loads fontspec
  \defaultfontfeatures{Scale=MatchLowercase}
  \defaultfontfeatures[\rmfamily]{Ligatures=TeX,Scale=1}
\fi
\usepackage{lmodern}
\ifPDFTeX\else
  % xetex/luatex font selection
\fi
% Use upquote if available, for straight quotes in verbatim environments
\IfFileExists{upquote.sty}{\usepackage{upquote}}{}
\IfFileExists{microtype.sty}{% use microtype if available
  \usepackage[]{microtype}
  \UseMicrotypeSet[protrusion]{basicmath} % disable protrusion for tt fonts
}{}
\makeatletter
\@ifundefined{KOMAClassName}{% if non-KOMA class
  \IfFileExists{parskip.sty}{%
    \usepackage{parskip}
  }{% else
    \setlength{\parindent}{0pt}
    \setlength{\parskip}{6pt plus 2pt minus 1pt}}
}{% if KOMA class
  \KOMAoptions{parskip=half}}
\makeatother
\usepackage{xcolor}
\usepackage[margin=25mm]{geometry}
\setlength{\emergencystretch}{3em} % prevent overfull lines
\providecommand{\tightlist}{%
  \setlength{\itemsep}{0pt}\setlength{\parskip}{0pt}}
\setcounter{secnumdepth}{-\maxdimen} % remove section numbering
\providecommand{\gt}{>}
\providecommand{\lt}{<}
\ifLuaTeX
  \usepackage{selnolig}  % disable illegal ligatures
\fi
\IfFileExists{bookmark.sty}{\usepackage{bookmark}}{\usepackage{hyperref}}
\IfFileExists{xurl.sty}{\usepackage{xurl}}{} % add URL line breaks if available
\urlstyle{same}
\hypersetup{
  hidelinks,
  pdfcreator={LaTeX via pandoc}}

\author{}
\date{}

\begin{document}

Linear and quadratic approximations

Reading route

Read from ``A tangent as a local model'' through the clock application,
trying A1 and A2 where they occur and A3 at the end of the application.
The later revisit A4 combines recall with method choice. All
\protect\hyperlink{hints}{hints} are together, followed by all
\protect\hyperlink{solutions}{solutions}; use the task links to return.
You may stop at a partial attempt and use its hint. The aim is to
construct a local approximation, keep a consistent order, and decide
what that approximation can justify.

This route uses familiar differentiation, limits, elementary algebra,
radians and integration. It assumes no earlier university lecture. A
basepoint is an input where the function and its derivatives are known
or easy to calculate. We will approximate nearby values by a polynomial
in the displacement from that point. The approximating polynomial has
exact coefficients; using it in place of the function generally gives an
estimate.

A tangent as a local model

Choose a fixed basepoint \(x_0\) and put \(h=x-x_0\). Thus \(x=x_0+h\);
the small quantity is \(h\), even if \(x_0\) is large. If \(f\) is
differentiable at \(x_0\), its tangent polynomial is

\[L(x)=f(x_0)+f'(x_0)(x-x_0).\]

Here \(f(x_0)\) is the height at contact, and \(f'(x_0)\) is the slope
there. Starting at the contact point, a horizontal displacement \(h\)
changes the line's height by \(f'(x_0)h\). Conversely, a line through
that point with that slope must have this equation. This is the meaning
of the lecture's figure \(y=b+a(x-x_0)\): \(b=f(x_0)\) and
\(a=f'(x_0)\). The curve and line share height and slope at contact;
they need not coincide elsewhere.

The derivative supplies a precise reason to approximate. Define the
error \(R_1(h)=f(x_0+h)-f(x_0)-f'(x_0)h\). For nonzero \(h\),

\[\frac{R_1(h)}h=\frac{f(x_0+h)-f(x_0)}h-f'(x_0)\longrightarrow0\quad(h\to0).\]

So the error becomes small compared with the displacement. We abbreviate
this property as \(R_1(h)=o(h)\): the symbol \(o(h)\) denotes an
unspecified remainder whose ratio to \(h\) tends to zero. It need not be
zero, and it is not a stated error bound at a particular nonzero input.
Writing \(f(x)\approx L(x)\) means using that local model. The guarantee
concerns approach to \(x_0\); it does not say the error must increase
monotonically at every step away from \(x_0\).

For \(f(x)=\ln x\) on \(x\gt 0\), choose \(x_0=1\). Since \(f(1)=0\) and
\(f'(1)=1\), the exact tangent polynomial is \(L(x)=x-1\). Therefore
\(\ln x\approx x-1\) near \(1\); for example, \(\ln(1.03)\approx0.03\).
To describe the same calculation with a zero basepoint, set \(u=x-1\).
Then \(x=1+u\), \(u\) is near zero, and \(\ln(1+u)\approx u\), with
\(u\gt -1\). The small input to this last formula is the change \(u\),
not the original input \(x\) near \(1\).

A1. First use

For f(x)=sqrt(x), use the tangent at basepoint x0=4 to estimate
sqrt(4.08). Give the whole tangent polynomial L(x), substitute the
target input, and explain what is exact and what is approximate. Use the
real square root on x\textgreater0. A sufficient response names the
basepoint and displacement and keeps the exact tangent equation distinct
from the estimate of f.

\protect\hyperlink{h1}{Hint A1} · \protect\hyperlink{s1}{Solution A1} ·
\protect\hyperlink{route}{Reading route}

Five basic local models

For the next formulas the basepoint is zero and \(x\) itself is the
displacement. Trigonometric inputs are in radians. The exponential and
logarithm inputs are dimensionless. The exponent \(r\) in the last
formula is fixed; for general real \(r\) we work near zero with
\(1+x\gt 0\).

For sine, \(f(0)=0\) and \(f'(0)=\cos0=1\), giving \(\sin x\approx x\).
For cosine, \(f(0)=1\) and \(f'(0)=-\sin0=0\), giving
\(\cos x\approx1\). The two tangent pictures mean different things: sine
has a sloping tangent through the origin, while cosine has a horizontal
tangent through height \(1\). A ``linear approximation'' includes a
constant polynomial when the slope is zero.

For \(e^x\), value and slope at zero are both \(1\), so
\(e^x\approx1+x\). For \(\ln(1+x)\), value zero and derivative
\(1/(1+x)\) give \(\ln(1+x)\approx x\). For \((1+x)^r\), value \(1\) and
derivative \(r(1+x)^{r-1}\) give \((1+x)^r\approx1+rx\). All five are
approximations near zero, not identities. The derivative argument above
gives an \(o(x)\) remainder in each case.

Substitution means replacing the entire small input. For instance,
\(e^u\approx1+u\) with \(u=-2x\) gives \(e^{-2x}\approx1-2x\), requiring
\(|-2x|\) small. Similarly \((1+x)^{-1/2}\approx1-x/2\). The identity
\(1/\sqrt{1+x}=(1+x)^{-1/2}\) requires \(x\gt -1\) and gives the
lecture's reciprocal-square-root shortcut directly. Alternatively, first
approximate \(\sqrt{1+x}\) by \(1+x/2\), then use
\((1+u)^{-1}\approx1-u\) with \(u=x/2\) for its reciprocal. Neither step
is an exact equality to the original function.

Combining approximations without inventing extra accuracy

The lecture's example is \(f(x)=e^{-2x}/\sqrt{1+x}\) for \(x\gt -1\).
Write it exactly as \(e^{-2x}(1+x)^{-1/2}\) and use the two linear
models:

\[f(x)\approx(1-2x)(1-x/2)=1-\frac52x+x^2.\]

The displayed multiplication is exact for the two polynomials, but its
\(x^2\) coefficient need not be the function's quadratic coefficient:
the factors' own quadratic contributions have not yet been included.
Keeping only constant and linear terms gives

\[L(x)=1-\frac52x,\qquad f(x)=1-\frac52x+o(x).\]

To justify combining the remainders, write the factors as \(1-2x+r(x)\)
and \(1-x/2+s(x)\), where \(r(x)/x\to0\) and \(s(x)/x\to0\). On
multiplying, the additional pieces are \(x^2\), \((1-2x)s(x)\),
\((1-x/2)r(x)\) and \(r(x)s(x)\). After division by \(x\) each tends to
zero: for example \(x^2/x=x\to0\) and \(r(x)s(x)/x=(r(x)/x)s(x)\to0\).
This licenses the stated first-order remainder. It also recovers the
exact facts \(f(0)=1\) and \(f'(0)=-5/2\) by continuity and the
derivative definition. A merely rough numerical fit would not justify
reading off an exact derivative this way.

Limits: cancellation makes the error matter

For the lecture's \(F(x)=(1+2x)^{10}\), the derivative definition
immediately gives

\[\lim_{x\to0}\frac{(1+2x)^{10}-1}{x}=F'(0)=20.\]

Indeed \(F(0)=1\) and \(F'(x)=20(1+2x)^9\). The approximation route
reaches the same answer: put \(u=2x\), \(r=10\) into the power model,
obtaining \(F(x)=1+20x+o(x)\). Subtracting \(1\) and dividing by nonzero
\(x\) leaves \(20+o(x)/x\to20\). This shows why an error statement
matters. Cancelling the constant does not turn \(\approx\) into an
equality; the remaining error must still vanish after the division used
in the limit.

Matching bending as well as slope

A linear model misses bending. Write a candidate polynomial in the same
displacement \(h=x-x_0\) as \(A+Bh+Ch^2\). Its value, first derivative
and second derivative at \(h=0\) are \(A,B,2C\). Matching these to
\(f(x_0),f'(x_0),f''(x_0)\) therefore gives

\[Q(x)=f(x_0)+f'(x_0)(x-x_0)+\frac{f''(x_0)}2(x-x_0)^2.\]

The coefficient is \(f''(x_0)/2\), because differentiating \(Ch^2\)
twice gives \(2C\). The whole displacement is squared. In particular, if
\(f(x)=a+bx+cx^2\) and \(x_0=0\), this construction gives \(a+bx+cx^2\)
exactly. At any basepoint it reproduces a polynomial of degree at most
two exactly, because that polynomial is already of the candidate form
when expressed in \(h\).

For a function whose first two derivatives are continuous near \(x_0\),
the error \(R_2(h)=f(x_0+h)-Q(x_0+h)\) satisfies \(R_2(h)/h^2\to0\). We
write this as \(R_2(h)=o(h^2)\). Here is why this stronger error claim
follows, beyond merely matching three numbers. The fundamental theorem
of calculus gives

\[f(x_0+h)-f(x_0)=\int_0^h f'(x_0+t)\,dt,\]

\[f'(x_0+t)=f'(x_0)+f''(x_0)t+\int_0^t\bigl[f''(x_0+s)-f''(x_0)\bigr]\,ds.\]

Insert the second equation into the first. Integrating the first two
terms gives \(f'(x_0)h+f''(x_0)h^2/2\), exactly the nonconstant terms of
\(Q\). The remainder is the double integral of the bracketed difference.
By continuity, for sufficiently small \(|h|\) the absolute value of that
difference is below any chosen positive number \(\varepsilon\)
throughout the interval between \(0\) and \(h\). The inner integral then
has magnitude at most \(\varepsilon|t|\), and the outer one at most
\(\varepsilon h^2/2\), for either sign of \(h\). Thus
\(|R_2(h)/h^2|\le\varepsilon/2\) as near zero as required, which proves
the claim. This establishes local accuracy without asserting a
particular numerical tolerance at a chosen input.

For cosine at zero, \(f(0)=1\), \(f'(0)=0\) and \(f''(0)=-1\). Hence
\(Q(x)=1-x^2/2\) and \(\cos x\approx1-x^2/2\). The constant tangent has
no bending; the negative quadratic term bends downward symmetrically to
either side, as cosine does near its maximum. This is the relation shown
by the source's parabola-and-cosine figure. ``Best fit'' here means
matching value, slope and second derivative at the basepoint; it does
not mean minimising errors over an arbitrary interval. A quadratic model
is useful when bending or a cancellation makes the second-order term
consequential, as well as when a closer local value is wanted.

The five quadratic models

At zero, sine has second derivative \(-\sin0=0\), cosine has second
derivative \(-\cos0=-1\), and the exponential has second derivative
\(e^0=1\). The resulting formulas are

\[\sin x\approx x,\qquad \cos x\approx1-\frac{x^2}{2},\qquad e^x\approx1+x+\frac{x^2}{2}.\]

A quadratic approximation means retaining terms through degree two; its
quadratic coefficient is allowed to be zero. For \(\ln(1+x)\),
differentiating \(f'(x)=(1+x)^{-1}\) gives \(f''(x)=-(1+x)^{-2}\), so
\((f(0),f'(0),f''(0))=(0,1,-1)\) and

\[\ln(1+x)\approx x-\frac{x^2}{2}.\]

For fixed \(r\), differentiating \(r(1+x)^{r-1}\) gives
\(r(r-1)(1+x)^{r-2}\). Therefore

\[(1+x)^r\approx1+rx+\frac{r(r-1)}2x^2.\]

These formulas have the same local domains and radian convention as the
linear models, and each has an \(o(x^2)\) remainder. The interval where
an infinite series converges and the interval where a short truncation
meets a desired tolerance are different questions. Even when
\(|x|\lt 1\) permits a nonterminating binomial series, that condition
alone does not certify an accurate two-term numerical estimate.
Coefficients and the required tolerance matter as well as smallness of
the input.

Returning to the combined function

To improve \(f(x)=e^{-2x}(1+x)^{-1/2}\), expand each factor through
degree two before multiplying. Substitution into the exponential and
power formulas gives

\[e^{-2x}=1-2x+2x^2+o(x^2),\]

\[(1+x)^{-1/2}=1-\frac{x}{2}+\frac38x^2+o(x^2).\]

The quadratic coefficient \(3/8\) comes from \((-1/2)(-3/2)/2\). To
build the product through degree two, a constant times a quadratic
contributes, and a linear times a linear contributes too. Thus

\[Q(x)=1+\left(-2-\frac12\right)x+\left(2+(-2)(-1/2)+\frac38\right)x^2
=1-\frac52x+\frac{27}{8}x^2.\]

Terms of degree three and above, divided by \(x^2\), tend to zero. A
factor remainder \(o(x^2)\) multiplied by a polynomial that stays
bounded near zero still has ratio to \(x^2\) tending to zero. The
product of the two remainders has that property too. Consequently
\(f(x)=Q(x)+o(x^2)\), and the coefficient \(27/8\) is justified. The
earlier isolated \(x^2\) term supplied only the mixed contribution
\(1\); it omitted \(2+3/8\) from the factors themselves.

A2. Combining and cancelling

Let g(x)=e\^{}(3x)/sqrt(1-2x), on the real domain x\textless1/2.
Construct its linear and quadratic polynomials about x=0. Show which
products contribute to the coefficient of x\^{}2. Then evaluate
lim\_(x-\textgreater0) {[}g(x)-g(0)-g'(0)x{]}/x\^{}2, and explain why a
first-order approximation alone does not determine this limit. A
sufficient response gives the two polynomials, the cross-term
calculation, and a remainder-based warrant for the limit, with the
substitutions' smallness conditions.

\protect\hyperlink{h2}{Hint A2} · \protect\hyperlink{s2}{Solution A2} ·
\protect\hyperlink{route}{Reading route}

A small parameter inside a physical model

The lecture's ``Planet Quirk'' satellite picture shows motion relative
to an observer; the direction arrow identifies relative motion, not an
additional term in the calculation. To give the clock comparison a
precise scope, use a constant-speed model without gravity or
acceleration. The laboratory is an inertial frame: a system of positions
and synchronised clocks in which an unforced body moves uniformly along
a straight line. Let \(T\gt 0\) be the time between two ticks recorded
by the travelling clock itself, and let \(T'\) be the interval assigned
to those same events by the laboratory. This is a comparison of event
times in the frame, not the arrival-time spacing of light signals at a
wristwatch. The special-relativistic law supplied for this example is

\[T'=\frac{T}{\sqrt{1-v^2/c^2}},\qquad c>0,\quad0\le v<c.\]

Here \(v\) is speed relative to the laboratory and \(c\) is light speed.
This physical law is an introduced premise, not a consequence of
differentiation. The inertial, moving-clock interpretation is consistent
with
\href{https://www.einstein-online.info/en/spotlight/light-clocks-time-dilation/}{Markus
Pössel's explanation of time dilation}. The mathematical work below
approximates that law. It is not a complete model of an orbiting clock
on a gravitating planet.

Define \(z=v^2/c^2\). Since both speeds use the same units, \(z\) and
\(T'/T\) are dimensionless. For \(v\ll c\), \(z\ll1\) and the small
input in the power formula is \(u=-z\), with exponent \(r=-1/2\).
Keeping through degree two in \(z\) gives

\[\frac{T'}T=(1-z)^{-1/2}\approx1+\frac z2+\frac38z^2,\]

\[T'-T\approx T\left(\frac{v^2}{2c^2}+\frac{3v^4}{8c^4}\right).\]

The sign of the leading change is positive: the laboratory assigns a
longer interval to those two ticks than the moving clock records. At
\(v=0\) the difference vanishes. ``Linear'' in \(z\) here already means
quadratic in \(v/c\); ``quadratic'' in \(z\) means quartic in \(v/c\).
This distinction prevents dropping the leading effect by looking for a
term proportional to \(v\).

With the lecture's rounded values \(v=4\ \mathrm{km/s}\) and
\(c=3\times10^5\ \mathrm{km/s}\),

\[z=\left(\frac4{300000}\right)^2\approx1.78\times10^{-10},\qquad \frac z2\approx8.89\times10^{-11},\qquad \frac38z^2\approx1.19\times10^{-20}.\]

The last two numbers are contributions to the ratio \(T'/T\).
Multiplying by the actual interval \(T\) gives time contributions; for
example the leading contribution for \(T=1\ \mathrm{s}\) is about
\(8.89\times10^{-11}\ \mathrm{s}\). Its smallness relative to \(1\) does
not by itself say whether a clock application needs it. Similarly, the
next term's size suggests the scale of the truncation error but does not
itself bound the sum of everything omitted. A required accuracy must be
compared with an actual error bound or a sufficiently justified error
estimate, in matching units.

Real GPS timing requires motion and gravitational effects.
\href{https://www.nist.gov/atomic-clocks/a-powerful-tool-for-science/putting-einstein-test}{NIST's
account} explains that they act in opposite directions for GPS
satellites and that the system corrects relativistic timing effects.
That supports the lecture's practical motivation while limiting what the
single displayed speed formula can establish. The lecture's 2006 comment
about the best atomic clocks is historical context, not a present-day
measurement limit used here. We need no such performance claim to
compare terms in the model.

A3. A changed clock comparison

Use the declared constant-speed, gravity-free model
T'=T/sqrt(1-(v/c)\^{}2), where T is the interval recorded by a clock
travelling with the object, T' is the interval assigned to the same two
ticks by the inertial laboratory, c\textgreater0, and
0\textless=v\textless c.~Set v/c=0.02 and T=1000 s. Estimate T'-T
retaining terms through (v/c)\^{}4. Find the size, in seconds, of the
quartic correction that is lost if only the leading nonzero correction
is kept. Compare it with an allowable absolute error of 10\^{}(-4) s.
Does the size of that one correction alone prove that the total error is
below the allowance? Explain. A sufficient response distinguishes the
dimensionless expansion variable, the clock-time difference, the order
in that variable versus the order in v/c, and an estimate of error from
a proved bound. No gravitational model or current clock-performance fact
is needed.

\protect\hyperlink{h3}{Hint A3} · \protect\hyperlink{s3}{Solution A3} ·
\protect\hyperlink{route}{Reading route}

Returning after a break

Try A4 at a later sitting, perhaps the next day; that timing is an
adjustable study suggestion. Reconstructing a formula checks retrieval
of what you studied. Selecting enough terms after a cancellation checks
transfer to the expression in front of you. If you need help, use the
hint, then return to the original task and complete the part you can
justify.

A4. Later retrieval and method choice

At a later sitting, first try without rereading: reconstruct the linear
and quadratic approximation formulas about a general basepoint a for a
function whose first two derivatives are continuous near a. Explain the
factor 1/2 by differentiating a quadratic polynomial. Then evaluate
lim\_(x-\textgreater0) {[}cos(x)-1{]}/x\^{}2 and explain why the
constant-plus-linear approximation alone cannot settle it. Finally state
whether going from linear to quadratic approximation changes the
polynomial for sin(x) about zero, and justify your answer with
derivatives at zero. Angles are in radians. A sufficient response
supplies the formulas, the coefficient-matching reason, the limit with
its remainder warrant, and the sine comparison. Partial attempts are
useful; mark the last step you can justify before opening help.

\protect\hyperlink{h4}{Hint A4} · \protect\hyperlink{s4}{Solution A4} ·
\protect\hyperlink{route}{Reading route}

Source and task provenance

This is an independently written teaching route based on the complete
seven-page
\href{https://ocw.mit.edu/courses/18-01-single-variable-calculus-fall-2006/40cb41807e6d67373f3299ce87c4d9b3_lec9.pdf}{MIT
OpenCourseWare 18.01, Fall 2006, Lecture 9: Linear and Quadratic
Approximations}, including its attribution page, four figures and
footnote. The logarithm, exponential/reciprocal-root product, power
limit, and Planet Quirk examples come from that source. Its attribution
and use information is at \href{https://ocw.mit.edu/terms/}{MIT OCW
Terms}. A1--A4 are generated practice, not quoted exam questions; the
clock task deliberately uses a changed speed and duration. The remainder
explanations and intermediate reasoning are supplied here to justify
approximation algebra and limits.

Two source formulas need care: some of its derivations print an equality
where only approximation is meant, and its first clock expansion drops
the multiplier \(T\). Here equality is reserved for exact expressions
(including expressions with remainders), and the dimensionless clock
expansion is explicitly for \(T'/T\). The source's printed exponent
\(4\) in the bracket explaining the clock binomial substitution is also
a typo; the stated exponent for that calculation is \(r=-1/2\). These
corrections preserve the mathematical operation without adopting the
printing errors.

Hints --- all tasks

Each hint advances one useful construction. Complete solutions are in
the separate group below.

Hint A1

Start with \(h=4.08-4\) and the values \(f(4)=2\),
\(f'(4)=1/(2\sqrt4)\). The line starts at height \(2\) and changes by
its slope times \(h\). The line's value is exact for that line; identify
which replacement makes the answer an approximation.

\protect\hyperlink{a1}{Return to A1} · \protect\hyperlink{s1}{Solution
A1}

Hint A2

Rewrite the denominator factor as \((1-2x)^{-1/2}\). Its power-formula
input is \(u=-2x\), while the exponential's input is \(u=3x\). For the
quadratic product, collect constant-times-quadratic, linear-times-linear
and quadratic-times-constant separately. Only after collecting those
terms subtract the value and slope terms in the limit.
\protect\hyperlink{quadratic}{The remainder explanation} tells you what
survives division by \(x^2\).

\protect\hyperlink{a2}{Return to A2} · \protect\hyperlink{s2}{Solution
A2}

Hint A3

First compute \(z=(0.02)^2\). The excess ratio is
\(T'/T-1\approx z/2+3z^2/8\); multiply each contribution by
\(1000\ \mathrm{s}\) before comparing with a tolerance in seconds. Keep
``the next term'' distinct from ``the whole error''.

\protect\hyperlink{a3}{Return to A3} · \protect\hyperlink{s3}{Solution
A3}

Hint A4

Use \(h=x-a\) and start from \(A+Bh+Ch^2\). Match its first two
derivatives and value at \(h=0\). In the cosine limit, subtraction
removes the value at zero, so the first surviving term matters. For
sine, inspect the second derivative at zero before deciding whether
there is a degree-two term.

\protect\hyperlink{a4}{Return to A4} · \protect\hyperlink{s4}{Solution
A4}

Complete reasoned solutions --- all tasks

These are model solutions to the generated tasks, not reports of
anyone's performance.

Solution A1

At \(x_0=4\), \(f(4)=2\) and \(f'(4)=1/4\), from \(f'(x)=1/(2\sqrt{x})\)
on \(x\gt 0\). Therefore the exact tangent polynomial is

\[L(x)=2+\frac14(x-4).\]

The displacement for the requested input is \(h=0.08\). Thus
\(L(4.08)=2+0.08/4=2.02\) exactly, and \(\sqrt{4.08}\approx2.02\) by the
local tangent approximation. The basepoint lies in the differentiable
domain and the displacement is small. Equality to the square root would
be false: \(2.02^2=4.0804\), rather than \(4.08\). This check
distinguishes exact line arithmetic from approximate use of the line.

\protect\hyperlink{a1}{Return to A1} · \protect\hyperlink{h1}{Hint A1} ·
\protect\hyperlink{route}{Reading route}

Solution A2

On \(x\lt 1/2\) the real denominator is positive, so
\(g(x)=e^{3x}(1-2x)^{-1/2}\). Near zero, both \(|3x|\) and \(|-2x|\) are
small. The substitutions into the established quadratic formulas give

\[e^{3x}=1+3x+\frac92x^2+o(x^2),\]

\[(1-2x)^{-1/2}=1+x+\frac32x^2+o(x^2).\]

For the second factor, the linear coefficient is \((-1/2)(-2)=1\) and
the quadratic coefficient is \([(-1/2)(-3/2)/2](-2)^2=3/2\). On
multiplication, the linear terms contribute \(3+1=4\), and the quadratic
terms contribute

\[\frac92+(3)(1)+\frac32=9.\]

Consequently \(L(x)=1+4x\) and \(Q(x)=1+4x+9x^2\). The established
product remainder rule gives \(g(x)=1+4x+9x^2+o(x^2)\). In particular
\(g(0)=1\) and \(g'(0)=4\): the first is direct substitution, and the
second follows by subtracting \(g(0)\), dividing by \(x\) and taking the
derivative limit. The requested limit is therefore

\[\lim_{x\to0}\frac{g(x)-g(0)-g'(0)x}{x^2}
=\lim_{x\to0}\left(9+\frac{o(x^2)}{x^2}\right)=9.\]

A first-order statement gives only \(g(x)=1+4x+o(x)\). Dividing its
unspecified error by \(x^2\) need not give zero or any particular finite
value; for example any term \(kx^2\) is \(o(x)\) but becomes \(k\) after
division by \(x^2\). The quadratic information is what determines this
limit. Differentiating the original product gives the independent check
\(g''(0)=18\), consistent with coefficient \(18/2=9\).

\protect\hyperlink{a2}{Return to A2} · \protect\hyperlink{h2}{Hint A2} ·
\protect\hyperlink{route}{Reading route}

Solution A3

The dimensionless small parameter is \(z=(v/c)^2=(0.02)^2=0.0004\).
Through degree two in \(z\), or degree four in \(v/c\),

\[T'-T\approx1000\left(\frac{0.0004}{2}+\frac38(0.0004)^2\right)\mathrm{s}
=0.20006000\ \mathrm{s}.\]

The leading term is \(0.2000\ \mathrm{s}\). The quartic correction
discarded by keeping only that term is
\(0.00006000\ \mathrm{s}=6.0\times10^{-5}\ \mathrm{s}\), which is
smaller than the allowance \(10^{-4}\ \mathrm{s}\). This comparison
suggests that keeping only the leading term may be adequate here, but
that single correction does not prove a bound on the total omitted
contribution. The local \(o(z^2)\) statement says how the remainder
behaves as \(z\to0\), not a numerical bound at \(z=0.0004\).
Certification would need a bound on the remaining terms or a controlled
comparison with the exact model. The requested estimate and the logical
limitation are both part of the answer; no claim about real clock
capabilities follows.

\protect\hyperlink{a3}{Return to A3} · \protect\hyperlink{h3}{Hint A3} ·
\protect\hyperlink{route}{Reading route}

Solution A4

With \(h=x-a\), the required polynomials are

\[L(x)=f(a)+f'(a)(x-a),\]

\[Q(x)=f(a)+f'(a)(x-a)+\frac{f''(a)}2(x-a)^2.\]

For \(A+Bh+Ch^2\), derivatives are \(B+2Ch\) and \(2C\). At \(h=0\),
matching the function, slope and second derivative sets \(A=f(a)\),
\(B=f'(a)\) and \(2C=f''(a)\). This explains the factor \(1/2\). Under
the stated continuity conditions the earlier integration argument
supplies the \(o(h^2)\) remainder.

For cosine at zero, the value and first two derivatives are \(1,0,-1\).
Thus \(\cos x=1-x^2/2+o(x^2)\). Subtract \(1\) and divide by \(x^2\) to
obtain

\[\lim_{x\to0}\frac{\cos x-1}{x^2}=-\frac12.\]

The linear model is the constant \(1\); it supplies no coefficient for
the first surviving term after the subtraction. Its error is known to be
\(o(x)\), which alone is insufficient after division by \(x^2\), for the
same reason demonstrated in Solution A2. For sine the value, slope and
second derivative at zero are \(0,1,0\). Therefore both its linear and
quadratic polynomials are \(x\). This says the degree-two coefficient
vanishes, not that \(\sin x=x\) identically or that there is no
higher-order error.

\protect\hyperlink{a4}{Return to A4} · \protect\hyperlink{h4}{Hint A4} ·
\protect\hyperlink{route}{Reading route}

\end{document}
